%COMANDO CHE DETERMINA IL TIPO DI DOCUMENTO CHE SI VUOLE CREARE

\documentclass[a4paper,12pt]{report}

%ELENCO DEI PACCHETTI UTILI PER LA STESURA DEL DOCUMENTO

\usepackage[utf8]{inputenc}
\usepackage[english, italian]{babel}
\usepackage{graphicx}
\usepackage{float}
\usepackage{tabularx}
\usepackage{makecell}
\usepackage{titlesec}
\usepackage{fancyhdr}
\usepackage{lastpage}
\usepackage{xurl}
\usepackage{hyperref}
\usepackage{geometry}
\usepackage[table,dvipsnames]{xcolor}
\setcounter{tocdepth}{5}
\setcounter{secnumdepth}{5}
\usepackage{caption}
\usepackage{etoolbox}% >= v2.1 2011-01-03
\usepackage{tikz}
\usepackage{listings}
\usepackage{scalerel}
\usepackage{longtable}
\usepackage{comment}
\usepackage{amssymb}
\usepackage{listings}

\titleformat{\chapter}[display]
{\normalfont\bfseries}{}{0pt}{\LARGE}

%COMANDO PER LA SPAZIATURA DEI TITOLI DAL BORDO DEL FOGLIO

\titlespacing*{\chapter}{0cm}{0cm}{0.2cm}

%COMANDO PER LA SPAZIATURA DEL TESTO DAI BORDI LATERALI

\geometry{
	left=20mm,
	right=20mm,
}

%COMADNO PER AVERE L'INDICE DEL NOME CHE SI VUOLE

\renewcommand{\contentsname}{Indice}

%COMADNI PER OTTENERE SUBSUBSECTION NUMERATE E PRESENTI NELL'INDICE

\setcounter{tocdepth}{5}
\setcounter{secnumdepth}{5}

%COMANDI PER OTTENERE HEADER E FOOTER

\pagestyle{plain}

\fancypagestyle{plain}{
	\fancyhf{}
	\lhead{\includegraphics[width=3cm]{../../immagini/minilogo.jpg}}
	\chead{}
	\rhead{\fontsize{12}{10}Verbale Esterno 2021-04-15}
	\lfoot{}
	\cfoot{\thepage\ di \pageref{LastPage}}
	\rfoot{}
}

\pagestyle{plain}

%COMANDI PER LINK

\hypersetup{
	colorlinks=true,
	linkcolor=black,
	filecolor=black,
	urlcolor=blue,
	citecolor=black,
}

%COMANDI PER SNIPPET DI CODICE

\BeforeBeginEnvironment{lstlisting}{\begin{mdframed}\vspace{-0.7em}}
	\AfterEndEnvironment{lstlisting}{\vspace{-0.5em}\end{mdframed}}

% needed for \lstcapt
\def\ifempty#1{\def\temparg{#1}\ifx\temparg\empty}

% make new caption command for listings
\usepackage{caption}
\newcommand{\lstcapt}[2][]{%
	\ifempty{#1}%
	\captionof{lstlisting}{#2}%
	\else%
	\captionof{lstlisting}[#1]{#2}%
	\fi%
	\vspace{0.75\baselineskip}%
}

\definecolor{atomlightorange}{rgb}{0.88,0.76,0.55}
\definecolor{atomdarkgrey}{RGB}{59,62,75}

% set listings
\lstset{%
	basicstyle=\footnotesize\ttfamily\color{atomlightorange},
	framesep=20pt,
	belowskip=10pt,
	aboveskip=10pt
}

% add frame environment
\usepackage[%
framemethod=tikz,
skipbelow=8pt,
skipabove=13pt
]{mdframed}
\mdfsetup{%
	leftmargin=0pt,
	rightmargin=0pt,
	backgroundcolor=atomdarkgrey,
	middlelinecolor=atomdarkgrey,
	roundcorner=6
}

%BIBLIOGRAFIA

\makeatletter
\def\thebibliography#1{\chapter*{Bibliografia\@mkboth
		{Bibliografia}{Bibliografia}}\list
	{[\arabic{enumi}]}{\settowidth\labelwidth{[#1]}\leftmargin\labelwidth
		\advance\leftmargin\labelsep
		\usecounter{enumi}}
	\def\newblock{\hskip .11em plus .33em minus .07em}
	\sloppy\clubpenalty4000\widowpenalty4000
	\sfcode`\.=1000\relax}
\makeatother

%DEFINIZIONE COLORI

\definecolor{airforceblue}{rgb}{0.36, 0.54, 0.66}


\begin{document}
	
	\makeatletter
	\begin{titlepage}
		\begin{center}
			\vspace*{-4cm}
			\author{Jawa Druids} 
			\title{Glossario}
			\date{} %LASCIARE QUESTO CAMPO VUOTO, SE LO TOLGO STAMPA LA DATA CORRENTE
			\includegraphics[width=0.5\linewidth]{../immagini/DRUIDSLOGO.jpg}\\[4ex]
			{\huge \bfseries  \@title }\\[2ex] 
			{\LARGE  \@author}\\[50ex]
			\vspace*{-9cm}
			\begin{table}[H]
				\renewcommand{\arraystretch}{1.4}
				\centering
				\begin{tabular}{r | l}
					\textbf{Versione} & v2.0.0 \\%RIGA PER INSERIRE LA VERSIONE ULTIMA DEL DOCUMENTO
					\textbf{Data approvazione} & 15-03-2021\\
					\textbf{Responsabile} & Mattia Cocco\\
					\textbf{Redattori} & \makecell[tl]{Mattia Cocco \\ Andrea Cecchin} \\
					\textbf{Verificatori} & \makecell[tl]{Andrea Dorigo } \\
					%MAKECELL SERVE PER POI ANDARE A CAPO ALL'INTERNO DELLA CELLA
					\textbf{Stato} & Approvato\\
					\textbf{Lista distribuzione} & \makecell[tl]{Jawa Druids \\ Prof. Tullio Vardanega \\ Prof. Riccardo Cardin \\ Sync Lab}\\
					\textbf{Uso} & Esterno          
				\end{tabular}
			\end{table}
			\vspace{0.1cm}
			\hfill \break
			\fontsize{17}{10}\textbf{Sommario} \\
			\vspace{0.1cm}
			Questo documento raccoglie le definizioni di alcuni termini utilizzati negli altri documenti al fine di risolvere possibili ambiguità ed incomprensioni.
		\end{center}
	\end{titlepage}
	\makeatother

\quad
\begin{center}
	\LARGE\textbf{Registro delle modifiche}
\end{center}

\def\tabularxcolumn#1{m{#1}}
{\rowcolors{2}{RawSienna!90!RawSienna!20}{RawSienna!70!RawSienna!40}


\begin{center}
	\renewcommand{\arraystretch}{1.4}
	\begin{longtable}[c]{|p{2cm-1\tabcolsep}|p{2cm}|p{3cm-2\tabcolsep}|p{2,5cm-2\tabcolsep}|p{2,5cm}|p{4cm-2\tabcolsep}|}
		\hline
		\rowcolor{airforceblue}
		\makecell[c]{\textbf{Versione}} & \makecell[c]{\textbf{Data}} & \makecell[c]{\textbf{Autore}} & \makecell[c]{\textbf{Ruolo}} & \makecell[c]{\textbf{Verificatore}} & \makecell[c]{\textbf{Modifica}}\\
		\hline
		\centering v1.0.0 & 07-02-2021 & Andrea Dorigo & \centering \textit{Analista} & \centering - & \textit{Approvazione del documento} \\
		\hline
		\centering v0.1.0 & 05-02-2021 & Emma Roveroni & \centering \textit{Analista} & Andrea Cecchin & \textit{Revisione complessiva del documento} \\
		\hline
		\centering v0.0.1 & 03-02-2021 & Emma Roveroni & \centering \textit{Analista} & Andrea Cecchin & \textit{Stesura del documento} \\
		\hline
		
	\end{longtable}
\end{center}

\tableofcontents{}

\chapter{A} \label{A}
\begin{itemize}
		\item \textbf{Applicazioni Web}\\
		Vedi Web-app$_G$.
		\item \textbf{Attività}\\
		Composizione di più compiti$_G$ semplici svolti da singoli, o gruppi, al fine di perseguire un obbiettivo comune.

	\end{itemize}

\chapter{B} \label{B}
\begin{itemize}
	\item \textbf{Back-End/Back end/Backend} \\
	Interfaccia con la quale il gestore di un sito web dinamico ne gestisce i contenuti e le funzionalità. A differenza del frontend$_G$, l'accesso al backend è riservato agli amministratori del sito che possono accedere dopo essersi autenticati.
	
	\item \textbf{Brainstorming} \\
	Metodo decisionale in cui, attraverso l’espressione libera delle proprie idee, mediante riunioni e confronti si ricerca una soluzione ad un dato problema.
	
	\item \textbf{Branch} \\
	Nel contesto informatico e di Git$_G$ si tratta di un ramo di lavoro.
	
		\item \textbf{Build automation} \\
	In informatica è l'atto di scrivere o automatizzare un'ampia varietà di compiti che gli sviluppatori software fanno nelle loro attività quotidiane di sviluppo.
\end{itemize}
\chapter{C} \label{C}
\begin{itemize}
	\item \textbf{Capitolato (d'appalto)}\\
	Documento tecnico, allegato generalmente ad un contratto d'appalto, a cui si fa riferimento per la definizione delle specifiche tecniche dei lavori che verranno successivamente svolti per effetto dello stesso contratto.
	
	\item \textbf{Committente}\\
	Soggetto che ordina ad altri l'esecuzione di un lavoro.
	
	\item \textbf{Compito}\\
	Rappresenta la componente elementare di un'attività, viene svolta da un soggetto per senso del dovere.
	Viene anche utilizzato al plurale: compiti.
\end{itemize}
\chapter{D} \label{D}
	\begin{itemize}		
		\item \textbf{Design Pattern}\\
		Soluzione progettuale generale ad un problema ricorrente. Si tratta di una descrizione, o modello logico, da applicare per la risoluzione di un problema che può presentarsi in diverse situazioni durante le fasi di progettazione e sviluppo del prodotto software.
		
		\item \textbf{Discord}\\
		Software gratuito utilizzato come canale di comunicazione da parte del gruppo di lavoro.
		Accessibile mediante smartphone, tablet e computer.
		Permette la comunicazione testuale, condivisione di file, chiamate vocali e video chiamate.
		Inoltre permette la creazione interna di canali personalizzati in modo da organizzare al meglio la comunicazione tra le persone.
		
		\item \textbf{Diagrammi di Gantt}\\
		Strumento di supporto alla gestione dei progetti, è costruito partendo da un asse orizzontale, a rappresentazione dell'arco temporale totale del progetto, suddiviso in fasi incrementali  e da un asse verticale, a rappresentazione delle mansioni o attività che costituiscono il progetto. Permette dunque la rappresentazione grafica di un calendario di attività$_G$, utile al fine di pianificare, coordinare e tracciare specifiche attività$_{\scaleto{G}{3pt}}$ in un progetto dando una chiara illustrazione dello stato d'avanzamento del progetto rappresentato.  Ad ogni attività$_{\scaleto{G}{3pt}}$ possono essere in generale associati una serie di attributi come la durata, le risorse ed il costo.
		
		\item \textbf{Drag}\\
		Parola inglese col significato di trascinare, nel contesto utilizzato indicato lo spostamento della heat map$_{\scaleto{G}{3pt}}$ negli assi orrizzontali e verticali.
	\end{itemize}
\chapter{E} \label{E}

\chapter{F} \label{F}
\begin{itemize}	
	\item \textbf{Fogli Google}\\
	Programma per fogli di calcolo facente parte della suite gratuita di Google Docs Editor offerta da Google.
		
	\item \textbf{Fornitore}\\
	Individuo, team o azienda incaricato di realizzare  il prodotto software richiesto dal proponente$_G$.		
	
	\item \textbf{Framework} \\
	Utilizzato per descrivere la struttura operativa nella quale viene elaborato un dato software.
	Un framework, in generale, include software di supporto, librerie, un linguaggio per gli script e altri software che possono aiutare a mettere insieme le varie componenti di un progetto.
	
	\item \textbf{Front-End/Front end/Frontend} \\
	Parte di un sistema software che gestisce l'interazione con l'utente o con sistemi esterni che producono dati di ingresso.
	Tali dati sono poi utilizzabili dal Back-end$_G$.


\end{itemize}
\chapter{G} \label{G}
\begin{itemize}
	\item \textbf{Git}\\
	Sistema di controllo gratuito a versione distribuita progettato per tenere traccia del lavoro svolto durante l'intero periodo di sviluppo del software.
	Utilizzato anche per tenere traccia di tutte le modifiche fatte nei file.
	I suoi punti di forza sono l'integrità dei dati e il supporto per flussi di lavoro distribuiti e non lineari.

	\item \textbf{Git-Flow}\\
	Si tratta di un ramo di Git$_G$, molto adatto alla collaborazione e scalarità del team di sviluppo.
	
	\item \textbf{Git-Hub}\\
	GitHub è un servizio di hosting per progetti software. Il nome deriva dal fatto che esso è una implementazione dello strumento di controllo versione distribuito Git$_G$.
	
	\item \textbf{Git-Kraken}\\
	Piattaforma grafica per Git$_G$.
	Utilizzata per tenere traccia graficamente della repository$_G$, vedere i branch$_G$, potendo anche scaricare il lavoro presente in locale, crearne di nuovi.
	Inoltre si può vedere lo storico di tutte le modifiche e la possibilità di fare nuovi commit.
	
	\item \textbf{GET}\\
	Il GET è una tipologia di richiesta al server che permette di ottenere informazioni riguardanti una risorsa.
	\end{itemize}
\chapter{H} \label{H}
\begin{itemize}
	\item \textbf{Heat Map} \\
	 Rappresentazione grafica dei dati dove i singoli valori contenuti in una matrice sono rappresentati da colori.
\end{itemize}
\chapter{I} \label{I}

\chapter{J} \label{J}
\begin{itemize}
	\item \textbf{Java} \\
	 Linguaggio di programmazione ad alto livello, orientato agli oggetti e a tipizzazione statica, che si appoggia sull'omonima piattaforma software di esecuzione, specificamente progettato per essere il più possibile indipendente dalla piattaforma hardware di esecuzione. Le principali caratteristiche di Java sono la portabilità, cioè il codice sorgente e' compilato in bytecode e puo' essere eseguito su ogni PC che ha JVM (Java Virtual Machine), e la robustezza. 
\end{itemize}
\chapter{K} \label{K}
\begin{itemize}
	\item \textbf{Kafka}\\
	Apache Kafka è una piattaforma open-source$_G$ di stream(si veda streaming$_G$) processing. Scritta in Java$_G$ e
	Scala e sviluppata dall’Apache Software Foundation, mira a creare una piattaforma a bassa
	latenza ed alta velocità per la gestione di feed dati in tempo reale.
	
\end{itemize}

\chapter{L} \label{L}
\begin{itemize}
	\item \textbf{Layout}\\
	In informatica si intende la disposizione degli elementi che costituiscono una pagina internet.
	
	\item \textbf{Leaflet} \\
	Libreria open-source$_G$ di JavaScript che permette la realizzazione di mappe interattive che funzionano efficientemente sia su desktop che su mobile.

	\item \textbf{Live-feed}\\
	Con questo termine si intendono delle immagini dal vivo trasmesse in diretta.

\end{itemize}

\chapter{M} \label{M}
\begin{itemize}
	\item \textbf{Machine Learning} \\
	Metodo di analisi dati che automatizza la costruzione di modelli analitici. È una branca dell'Intelligenza Artificiale e si basa sull'idea che i sistemi possono imparare dai dati, identificare modelli autonomamente e prendere decisioni con un intervento umano ridotto al minimo.
\end{itemize}
\chapter{N} \label{N}
	\begin{itemize}
		\item \textbf{Norma}\\
		Regola di condotta, stabilità da un singolo soggetto o da un collettivo, al fine di regolare un'attività pratica o di indicare il procedimento da seguire in determinate circostanze.
	\end{itemize}
\chapter{O} \label{O}
\begin{itemize}
	\item \textbf{Open-Source} \\
	Un software open-source è reso tale per mezzo di una licenza attraverso cui i detentori dei diritti favoriscono la modifica, lo studio, l'utilizzo e la redistribuzione del codice sorgente.
	
	\item \textbf{Outlier}\\
	Termine utilizzato in statistica per definire, in un insieme di osservazioni, un valore anomalo e aberrante. Un valore quindi chiaramente distante dalle altre osservazioni disponibili.
\end{itemize}
\chapter{P} \label{P}
	\begin{itemize}
		\item \textbf{Product Baseline}\\
		Tradotto: linea di base del prodotto.
		Contiene tutti i contenuti rilasciabili del progetto.
		
		\item \textbf{Proof of Concept}\\
		Prototipo di software che ha lo scopo di dimostrare che il progetto può essere sviluppato in modo conforme alle richieste.
		
		\item \textbf{Procedura}\\
		Nell'informatica s'intende una sequenza ordinata di operazioni da eseguire al fine di raggiungere un determinato scopo.

		\item \textbf{Proponente}\\
		Soggetto che propone il capitolato$_G$ d'appalto per un progetto.
		
		\item \textbf{Python}\\
		Python è un linguaggio di programmazione ad alto livello, rilasciato pubblicamente per la prima volta nel 1991 dal suo creatore Guido van Rossum, supporta diversi paradigmi di programmazione, come quello orientato agli oggetti (con supporto all'ereditarietà multipla), quello imperativo e quello funzionale, ed offre una tipizzazione dinamica forte.
		Python è un linguaggio pseudocompilato: un interprete si occupa di analizzare il codice sorgente e, se sintatticamente corretto, di eseguirlo. Questa caratteristica rende Python un linguaggio portabile. Una volta scritto un sorgente, esso può essere interpretato ed eseguito sulla gran parte delle piattaforme attualmente utilizzate, semplicemente basta la presenza della versione corretta dell’interprete.
		
		\item \textbf{Pop-up}\\
		Termine inglese che indica gli elementi dell'interfaccia grafica che compaiono automaticamente durante l'uso di una applicazione.
	\end{itemize}
\chapter{Q} \label{Q} 

\chapter{R} \label{R}
	\begin{itemize}
		\item \textbf{Repository} \\
		Ambiente di un sistema informativo in cui vengono conservati e gestiti file, documenti e metadati relativi ad un’attività$_G$ di progetto.
		
		\item \textbf{Requisito}\\
		Una condizione, o capacità, che deve essere verificata o posseduta da un sistema o un componente di esso per soddisfare un contratto, uno standard, una specifica o qualsiasi altro documento formalmente specificato. L'insieme di tutti i requisiti formano la base del successivo sviluppo del sistema o del componente.
		
	\end{itemize}
\chapter{S} \label{S}
	
\chapter{T} \label{T}
	\begin{itemize}
		\item \textbf{Trello}\\
		Sito internet che permette di organizzare i task per ciascun membro del gruppo di lavoro.
		Al suo interno è possibile creare delle schede contenenti il task da fare.
		Inoltre nelle schede è possibile inserire il nome della persona adibita a svolgere quel task, allegati, commenti ed etichette per differenziare la tipologia dell'attività inserita.
		Oltre a ciò, \textit{trello} dà la possibilità di poter spostare queste schede su varie colonne per poter facilitare ancora di più l'organizzazione a livello visivo.
	\end{itemize}
\chapter{U} \label{U}

\chapter{V} \label{V}

\chapter{W} \label{W}
	\begin{itemize}
		\item \textbf{wrapper}
	\end{itemize}
\chapter{X} \label{X}

\chapter{Y} \label{Y}

\chapter{Z} \label{Z}



%COMANDO PER LA CREAZIONE DELL'INDICE

%PER RENDERE PIÙ CHIARA LA STESURA DEI DOCUMENTI È MEGLIO LASCIARE SEPARATI IN FILE DIVERSI OGNI CAPITOLO

% \input{esempio} -- esempio di codice per inserire un nuovo capitolo

\end{document}
